\begin{frame}{Rdzeń}

Rdzeń grafu \(H\) to najmniejszy jego podgraf, do którego \(H\) ma homomorfizm.

\vspace{1em}
\begin{figure}
    \centering
    % \begin{tikzpicture}[
%     scale=0.85,
%     line width=1pt,
%     minimum size=0.06cm
% ]

% % pętla
% % =====================================================================
% \begin{scope}[shift={(0,0)}]

%     \node[circle, line width=1pt, draw=white, fill=black] (v1) at (1,1) {};

%     \draw[->] (v1) to [out=90, in=0, looseness=10] (v1);

% \end{scope}


% % cykl
% % =====================================================================
% \begin{scope}[shift={(5,0)}]

%     \node[circle, line width=1pt, draw=white, fill=black] (v1) at (0,2) {};
%     \node[circle, line width=1pt, draw=white, fill=black] (v2) at (2,2) {};
%     \node[circle, line width=1pt, draw=white, fill=black] (v3) at (2,0) {};
%     \node[circle, line width=1pt, draw=white, fill=black] (v4) at (0,0) {};

%     \draw[->] (v1) -- (v2);
%     \draw[->] (v2) -- (v3);
%     \draw[->] (v3) -- (v4);
%     \draw[->] (v4) -- (v1);

% \end{scope}


% % klika
% % =====================================================================
% \begin{scope}[shift={(10,0)}]

%     \node[circle, line width=1pt, draw=white, fill=black] (v1) at (0,2) {};
%     \node[circle, line width=1pt, draw=white, fill=black] (v2) at (2,2) {};
%     \node[circle, line width=1pt, draw=white, fill=black] (v3) at (2,0) {};
%     \node[circle, line width=1pt, draw=white, fill=black] (v4) at (0,0) {};

%     \draw[<->] (v1) -- (v2);
%     \draw[<->] (v2) -- (v3);
%     \draw[<->] (v3) -- (v4);
%     \draw[<->] (v4) -- (v1);

%     \draw[<->] (v1) -- (v3);
%     \draw[<->] (v2) -- (v4);

% \end{scope}

% \end{tikzpicture}


\begin{tikzpicture}[
    scale=1.1,
    line width=1pt,
    minimum size=0.08cm
]
% \usetikzlibrary{backgrounds}

\begin{scope}[shift={(6,0)}]

    \node[circle, draw=black, fill=cfyel!85]  (v1) at (0,2) {};
    \node[circle, draw=black, fill=cfblue!85] (v2) at (2,2) {};
    \node[circle, draw=black, fill=cfyel!85]  (v3) at (2,0) {};
    \node[circle, draw=black, fill=cfblue!85] (v4) at (0,0) {};
    \node[circle, draw=black, fill=cfred!85]  (v5) at (1,1) {};

    \draw[<-] (v1) -- (v2);
    \draw[->] (v2) -- (v3);
    \draw[<-] (v3) -- (v4);
    \draw[->] (v4) -- (v1);

    \draw[->] (v5) -- (v2);
    \draw[->] (v3) -- (v5);

    \node at (1,-1) {\(H\)};
    \node at (3.6,1) {\(\core(H)\)};

\end{scope}

% highlight
\begin{scope}[on background layer]
    \fill[cfgrey!40, rounded corners=20pt]
        ([shift={(-1,0)}] v5.center) --
        ([shift={(0.5,1)}] v2.center) --
        ([shift={(0.5,-1)}] v3.center) --
        cycle;
\end{scope}

\end{tikzpicture}
\end{figure}

\end{frame}

% --------------------------------------------------------------

\begin{frame}{Rdzeń}

Graf \(G\) jest \(H\)-kolorowalny wtw \(G\) jest \(\core(H)\)-kolorowalny.

\vspace{1em}
\begin{figure}
    \centering
    \begin{tikzpicture}[
    scale=0.9,
    line width=1pt,
    minimum size=0.08cm
]
% \usetikzlibrary{backgrounds}

% graf G
% =====================================================================
\begin{scope}[shift={(0,0)}]

    \node[circle, draw=black, fill=cfblue!85] (v1) at (-0.3,1.2) {};
    \node[circle, draw=black, fill=cfyel!85]  (v2) at (1,2.2) {};
    \node[circle, draw=black, fill=cfblue!85] (v3) at (2.3,1.2) {};
    \node[circle, draw=black, fill=cfyel!85]  (v4) at (1.7,-0.2) {};
    \node[circle, draw=black, fill=cfred!85]  (v5) at (0.3,-0.2) {};
    \node[circle, draw=black, fill=cfred!85]  (v6) at (1,0.9) {};

    \draw[->] (v1) -- (v2);
    \draw[<-] (v2) -- (v3);
    \draw[->] (v3) -- (v4);
    \draw[->] (v4) -- (v5);
    \draw[->] (v5) -- (v1);

    \draw[<-] (v1) -- (v6);
    \draw[->] (v4) -- (v6);

    \node at (1,-1) {\(G\)};

\end{scope}

% G -> H
\draw[->, line width=2pt, cfgrey] (3,1) to (5,1);
\node[gray] at (4,1.5) {hom};


% graf H
% =====================================================================
\begin{scope}[shift={(6,0)}]

    \node[circle, draw=black, fill=cfyel!85]  (v1) at (0,2) {};
    \node[circle, draw=black, fill=cfblue!85] (v2) at (2,2) {};
    \node[circle, draw=black, fill=cfyel!85]  (v3) at (2,0) {};
    \node[circle, draw=black, fill=cfblue!85] (v4) at (0,0) {};
    \node[circle, draw=black, fill=cfred!85]  (v5) at (1,1) {};

    \draw[<-] (v1) -- (v2);
    \draw[->] (v2) -- (v3);
    \draw[<-] (v3) -- (v4);
    \draw[->] (v4) -- (v1);

    \draw[->] (v5) -- (v2);
    \draw[->] (v3) -- (v5);

    \node at (1,-1) {\(H\)};

\end{scope}

% highlight
\begin{scope}[on background layer]
    \fill[cfgrey!40, rounded corners=20pt]
        ([shift={(-1,0)}] v5.center) --
        ([shift={(0.5,1)}] v2.center) --
        ([shift={(0.5,-1)}] v3.center) --
        cycle;
\end{scope}


% graf core(H)
% =====================================================================
\begin{scope}[shift={(11,0)}]

    \node[circle, draw=black, fill=cfblue!85] (v1) at (0,0.2) {};
    \node[circle, draw=black, fill=cfyel!85]  (v2) at (1,1.8) {};
    \node[circle, draw=black, fill=cfred!85]  (v3) at (2,0.2) {};

    \draw[->] (v1) -- (v2);
    \draw[->] (v2) -- (v3);
    \draw[->] (v3) -- (v1);

    \node at (1,-1) {\(\core(H)\)};

\end{scope}

% highlight
\begin{scope}[on background layer]
    \fill[cfgrey!40, rounded corners=20pt]
        ([shift={(-1,-0.5)}] v1.center) --
        ([shift={(0,1)}] v2.center) --
        ([shift={(1,-0.5)}] v3.center) --
        cycle;
\end{scope}


% H -> core(H)
\draw[->, line width=2pt, cfgrey] (8.6,3) to[bend left=25] (11,3);
\node[gray] at (9.8,3.7) {hom};

% core(H) -> H
\draw[<-, line width=2pt, cfgrey] (8.6,-1.5) to[bend right=25] (11,-1.5);
\node[gray] at (9.8,-2.15) {hom};

\end{tikzpicture}
\end{figure}

\vspace{1em}

\LARGE
\[ \underset{\text{trudny}}{G} \ \overset{\text{hom}}{\longrightarrow} \ H \ \overset{\text{hom}}{\longrightarrow} \ \underset{\text{trudny}}{\core(H)} \]

\end{frame}
